\documentclass[12pt, a4paper]{article}
\usepackage[UTF8]{ctex}
\usepackage{geometry}
\usepackage{hyperref}
\usepackage{amsmath, amssymb}
\usepackage{enumitem}
\usepackage{fancyhdr}

\geometry{margin=2.5cm}

\hypersetup{
    colorlinks=true,
    linkcolor=blue,
    urlcolor=blue
}

\pagestyle{fancy}
\fancyhf{}
\fancyhead[L]{算法题目整理}
\fancyhead[R]{\thepage}
\renewcommand{\headrulewidth}{0.4pt}

\newcommand{\category}[1]{%
  \vspace{1.2cm}
  {\noindent\Large\bfseries #1}
  \vspace{0.2em}\par
  \noindent\rule{\textwidth}{0.8pt}
  \vspace{0.2em}
}

\title{\textbf{算法题目整理}}
\author{}
\date{}

\begin{document}
\maketitle
\thispagestyle{empty}


% ╔════════════════════════════════════════════════════════════╗
\category{一、树 / 二叉树}
% ╚════════════════════════════════════════════════════════════╝

\section{极北之境：奥术协议的坍缩（万灵之楔）}

\noindent\textbf{题目：}

在永恒冻土的尽头，屹立着一座名为"逻辑圣所"的塔。塔顶悬浮着一台古老的算力机——万灵之楔。
作为最后的秘法翻译官，你手中握着 $n$ 颗含有不同以太浓度的"源质核心"（\texttt{gem[i]}）。为了重启圣所，你必须将这些核心嵌入一个完全二叉树结构的插槽组中。

\textbf{协议约束 (Structure Protocol)}

\textbf{节点极性}：圣所提供无限的"中转节点"。每个中转节点必须消耗能量产生两个分支（左、右子节点），否则能量将因不平衡而爆炸。

\textbf{源质嵌入}：所有的 $n$ 颗源质核心必须全部且只能作为该树的叶子节点。

\textbf{序列坍缩 (The Sequence Collapse)：}

能量流从根节点注入，遵循"左弦优先"的遍历法则。每当能量状态发生迁移，圣所的示波器会记录下一组脉冲信号：
\begin{itemize}[leftmargin=2em]
  \item \textbf{注入脉冲}：能量进入任何节点（中转或核心）时，记录数字 1。
  \item \textbf{共鸣脉冲}：能量触碰核心（叶子）时，立即记录该核心的以太浓度值 \texttt{gem[i]}。
  \item \textbf{回归脉冲}：能量离开某个节点并回溯至父节点时，记录数字 9。
\end{itemize}

\textbf{最终裁定 (The Verdict)}

所有脉冲信号按记录顺序首尾相接，将坍缩为一个庞大的十进制整数 $V$。圣所的防御机制由两个神秘常数 $P$（模数）和 $T$（目标残差）控制。

\textbf{任务}：计算有多少种不同的核心排列与树形拓扑结构，能够产生满足 $V \equiv T \pmod{P}$ 的重启序列。

\textbf{注意：}
\begin{itemize}[leftmargin=2em]
  \item 即便两颗核心的数值相同，只要它们的物理编号不同，就视为不同的组合。
  \item 树的拓扑结构不同，即视为不同方案。
\end{itemize}
示例：

输入：gem = [2,3] p = 100000007 

t = 11391299

输出：1



\section{文件系统的备份日志}

\noindent\textbf{题目：}

在一个古老的服务器中，存在着一套文件目录系统，以树形结构组织所有的文件夹。系统管理员整理了各文件夹之间的层级关系，\texttt{parents[i]} 表示编号为 $i$ 的文件夹的父文件夹编号（根目录的父文件夹为 $-1$）。

为了备份这套目录结构，管理员设计了一种极简的日志格式，可以用一个 01 字符串来记录整个目录树的构建过程，规则如下：

初始时只有一个根目录，光标指向根目录。

依次读取字符串中的每个字符：
\begin{itemize}[leftmargin=2em]
  \item 如果是 \texttt{0}，则在当前文件夹下创建一个子文件夹，并将光标移动到这个新创建的子文件夹
  \item 如果是 \texttt{1}，则将光标返回到当前文件夹的父文件夹
\end{itemize}

现在需要用上述日志格式，完整地备份这套目录结构。请返回能够完成备份的所有可行字符串中，字典序最小的那个 01 字符串。

注意：备份结束时，光标可以停留在任意文件夹，不要求必须回到根目录。

示例：

输入：parents = [-1,0,0,2]

输出："00110"



\section{二叉树删除重组后的最大层和}

\noindent\textbf{题目：}

给定一棵二叉树的层序遍历序列 \texttt{root}，其中 \texttt{null} 表示空节点，每个节点对应一个整数值。

你可以对该二叉树执行最多一次"删除并重组"的操作，规则如下：

\textbf{删除限制}：只能选择删除叶子节点或只有一个子节点的节点。

\textbf{结构调整}：如果被选中的节点有一个子节点，则该子节点将提升到被删除节点的位置，成为其父节点的子节点（即以该子节点为根的子树上移）。

请问：在进行一次上述操作或选择不操作后，整棵二叉树的「层和」最大是多少？

\textbf{定义：}「层和」：二叉树中同一深度（层）上所有节点值的总和。

\textbf{示例 1：}

输入：\texttt{root = [6,0,3,null,8]}

输出：\texttt{11}

解释：删除值为 0 的节点（其右子节点 8 提升上来），此时第二层（深度为1）的节点值之和为 $3 + 8 = 11$，是所有层中的最大值。

\textbf{示例 2：}

输入：\texttt{root = [5,6,2,4,null,null,1,3,5]}

输出：\texttt{9}

\textbf{示例 3：}

输入：\texttt{root = [-5,1,7]}

输出：\texttt{8}

\textbf{提示：}
\begin{itemize}[leftmargin=2em]
  \item $2 \leq$ 树中节点个数 $\leq 10^5$
  \item $-10000 \leq$ 树中节点的值 $\leq 10000$
\end{itemize}



\section{企业并购结构设计}

\noindent\textbf{题目：}

\textbf{题目描述}

某投资公司正在设计一个企业并购结构。已知公司拥有足够数量的控股公司和 $n$ 个子公司，\texttt{gem[i]} 表示第 $i$ 个子公司的估值。现在需要使用这些控股公司和子公司组合成一个并购层级结构 (二叉树)，其组装规则为：
\begin{itemize}[leftmargin=2em]
  \item 控股公司将作为层级结构中的非叶子节点，且每个控股公司必须拥有 2 个下属实体；
  \item 子公司将作为层级结构中的叶子节点，所有的子公司都必须被纳入。
\end{itemize}

审计追踪从根节点开始，而后将按如下规则进行记录：
\begin{itemize}[leftmargin=2em]
  \item 若审计首次到达该节点时：
  \item 记录数字 1；
  \item 若该节点为叶节点 (子公司)，将额外记录该子公司的估值；
  \item 若存在未审计的下属实体，则选取未审计的一个实体（优先选取左侧实体）进入；
  \item 若无下属实体或所有下属实体均已审计过，此时记录 9，并回到当前节点的上级节点（若存在）。
\end{itemize}

如果最终记下的数依序连接成一个整数 \texttt{num}，满足 $\texttt{num} \bmod p = \texttt{target}$，则视为符合税务代码合规要求。请问有多少种并购结构的组装方案，可以使得最终记录下的数字符合税务要求。

\textbf{注意：}
\begin{itemize}[leftmargin=2em]
  \item 两个结构不同的并购层级，作为不同的组装方案
  \item 两个结构相同的并购层级且存在某个相同位置处的子公司编号不同，也作为不同的组装方案
  \item 可能存在估值相同的两个子公司
\end{itemize}

\textbf{示例 1:}

输入：\texttt{gem = [2,3], p = 100000007, target = 11391299}

输出：\texttt{1}

解释：包含 2 个子公司的并购结构只有一种。假设 B、C 子公司的估值分别为 3、2，审计追踪结果对应 target 为 11391299，$11391299 \% 100000007 = 11391299$，满足税务要求; 假设 B、C 子公司的估值分别为 2、3，审计追踪结果对应 target 为 11291399; $11291399 \% 100000007 = 11291399$，不满足税务要求；因此只存在 1 种方案，返回 1。

\textbf{示例 2:}

输入：\texttt{gem = [3,21,3], p = 7, target = 5}

输出：\texttt{4}

\textbf{提示：}
\begin{itemize}[leftmargin=2em]
  \item $1 \leq \texttt{gem.length} \leq 9$
  \item $0 \leq \texttt{gem[i]} \leq 10^9$
  \item $1 \leq p \leq 10^9$，保证 $p$ 为素数。
  \item $0 \leq \texttt{target} < p$
  \item 存在 2 组 \texttt{gem.length} $== 9$ 的用例
\end{itemize}



\section{仿生神经网络最大激活路径}

\noindent\textbf{题目：}

科研团队设计了一种仿生神经网络，其拓扑结构为一棵二叉树。每个神经元（节点）有一个激活值（整数，可以为负数——表示该神经元处于抑制状态）。

定义一条「激活路径」为树中任意两个神经元之间的简单路径（也可以只包含一个神经元），路径上的每个神经元最多经过一次。路径的激活总值为路径上所有神经元激活值之和。

请找出该神经网络中激活总值最大的路径，返回其激活总值。

\textbf{示例 1：}

输入：\texttt{neuralNet = [1, 2, 3]}（层序遍历表示）

输出：\texttt{6}

解释：最优激活路径为 $2 \to 1 \to 3$，总值 $= 2+1+3 = 6$。

\textbf{示例 2：}

输入：\texttt{neuralNet = [-10, 9, 20, null, null, 15, 7]}

输出：\texttt{42}

解释：最优激活路径为 $15 \to 20 \to 7$，总值 $= 15+20+7 = 42$。

\textbf{提示：}
\begin{itemize}[leftmargin=2em]
  \item 节点数量范围 $[1, 3 \times 10^4]$
  \item $-1000 \leq$ 节点激活值 $\leq 1000$
\end{itemize}



% ╔════════════════════════════════════════════════════════════╗
\category{二、动态规划}
% ╚════════════════════════════════════════════════════════════╝

\section{混沌波动法阵}

\noindent\textbf{题目：}

"杰洛特，别再抱怨传送门了。如果你想甩掉艾瑞汀和他的红色骑士，这是唯一的办法。"

叶奈法正在着手构建一个极其复杂的连续时空跳跃法阵，试图掩盖希里的行踪。为了彻底迷惑狂猎的首席导航员卡兰希尔，这次传送必须违背常规的魔法流动规律。

这个法阵由 $n$ 个连续的传送节点组成。叶奈法需要为每一个节点注入混沌能量，能量值的范围必须精确地控制在 $[l, r]$ 之间。

为了防止狂猎的法师通过魔法残痕锁定坐标，杰洛特必须协助叶奈法确保能量序列满足以下"混沌波动法则"：

\begin{enumerate}[leftmargin=2em]
  \item \textbf{拒绝共振}：为了防止空间坍塌，任意两个相邻节点的能量值不能相同。
  \item \textbf{打破轨迹}：卡兰希尔擅长通过线性魔力残留预判下一个落点。因此，能量流绝不能呈现规律性的增强或减弱。即：任意连续的三个节点，其能量值不能构成"严格递增"（能量持续暴涨）或"严格递减"（能量持续衰退）的序列。
\end{enumerate}

换句话说，每一个中间节点的能量，必须比两边都高（波峰），或者比两边都低（波谷），以此形成剧烈的锯齿状波动来扰乱追踪。

\textbf{任务}：作为对此一窍不通但擅长思考的猎魔人，请你告诉叶奈法，在给定的能量区间 $[l, r]$ 内，总共有多少种符合"混沌波动法则"的传送门构建方案？

（由于方案数可能像史凯利格群岛的雨点一样多，请将结果对 $10^9 + 7$ 取余后回答。）

\textbf{示例映射：}

输入：\texttt{n = 3, l = 1, r = 3}

情境：需要建立 3 个传送门，能量可选 1, 2, 3。

解释：
\begin{itemize}[leftmargin=2em]
  \item $[1, 2, 3]$ 是禁止的（严格递增，卡兰希尔会瞬间算出你去向）。
  \item $[3, 2, 1]$ 是禁止的（严格递减）。
  \item $[1, 2, 1]$ 是允许的（能量由低到高，再突然降低，成功迷惑敌人）。
  \item $[2, 1, 3]$ 是允许的（能量由高到低，再突然升高）。
\end{itemize}

最终符合条件的方案共 10 种。



\section{买卖股票的最佳时机 IV}

\noindent\textbf{题目：}

考虑一个点集，其中第 $i$ 个点的坐标为 $(i, \texttt{nums}[i])$。选择 $k$ 个互不相交的线段，每个线段连接两个点，要求连接的两个点的 $x$ 坐标和 $y$ 坐标都满足后者大于前者。最大化所有线段 $y$ 坐标差值之和。

\textbf{示例：}

输入：\texttt{k = 2, nums = [2, 4, 1]}

输出：\texttt{2}

解释：第一条线 $[0,2] \to [1,4]$，result $= 4 - 2 = 2$



\section{通信网络频段分配}

\noindent\textbf{题目：}

在一个 $n \times m$ 的通信网格中，每个节点可以选择：

不发射（\texttt{.}）

使用频段 L

使用频段 R

若在同一行或同一列中：两个使用不同频段的节点中间恰好隔一个正在发射的节点，则会产生二阶干扰，系统不稳定。部分节点频段未知（\texttt{?}）。

问：一共有多少种频段分配方案，使系统稳定？

示例：输入：n = 3, m = 3, net = ["..L","..R","?L?"]

输出：5



\section{长电文密钥提取方案数}

\noindent\textbf{题目：}

情报部门截获了一段极长的加密电文 \texttt{intercept}（由英文字母组成）。密码分析师需要从电文中按原始出现顺序（不要求连续）选取若干字符，拼成一个已知的密钥短语 \texttt{key}。

不同的选取位置组合视为不同的提取方案，即使最终拼出的字符串内容相同。

请计算一共有多少种不同的提取方案可以从 \texttt{intercept} 中提取出 \texttt{key}。答案可能很大，请直接返回整数。

\textbf{示例 1：}

输入：\texttt{intercept = "rabbbit", key = "rabbit"}

输出：\texttt{3}

解释：有 3 种方式从 "rabbbit" 中提取 "rabbit"（选择第1/2/3个b中的不同两个）。

\textbf{示例 2：}

输入：\texttt{intercept = "babgbag", key = "bag"}

输出：\texttt{5}

\textbf{提示：}
\begin{itemize}[leftmargin=2em]
  \item $1 \leq |\texttt{intercept}|, |\texttt{key}| \leq 1000$
  \item 两者仅含英文字母
\end{itemize}



\section{双周期太阳帆能量采集}

\noindent\textbf{题目：}

一艘太阳帆飞船在环绕恒星飞行过程中，每个离散时刻记录了太阳辐射强度值，形成数组 \texttt{radiation}。飞船的储能系统有一个工程限制：在整个飞行过程中最多执行两次充放电周期。

每个充放电周期包含两个操作：
\begin{itemize}[leftmargin=2em]
  \item 开始充电：选择某个时刻 $i$，以当时的辐射强度 \texttt{radiation[$i$]} 作为充电基准。
  \item 释放储能：选择之后的某个时刻 $j > i$，收益为 $\texttt{radiation}[j] - \texttt{radiation}[i]$。
\end{itemize}

两个充放电周期不能时间重叠（即第一次释放必须在第二次充电之前或同时发生）。也可以选择只执行一次或零次周期。

请计算飞船通过最优策略能获取的最大总收益。

\textbf{示例 1：}

输入：\texttt{radiation = [3,3,5,0,0,3,1,4]}

输出：\texttt{6}

解释：第一次周期：时刻4充电(0)，时刻6释放(3)，收益3；第二次周期：时刻7充电(1)，时刻8释放(4)，收益3。总计6。

\textbf{示例 2：}

输入：\texttt{radiation = [1,2,3,4,5]}

输出：\texttt{4}

解释：一次周期即可：时刻1充电(1)，时刻5释放(5)，收益4。

\textbf{提示：}
\begin{itemize}[leftmargin=2em]
  \item $1 \leq |\texttt{radiation}| \leq 10^5$
  \item $0 \leq \texttt{radiation}[i] \leq 10^5$
\end{itemize}



\section{承重墙拆除能量最大化}

\noindent\textbf{题目：}

一栋老旧建筑的走廊中有 $n$ 面承重墙排成一排，每面墙有一个结构系数（正整数），存储在数组 \texttt{walls} 中。

拆迁队需要按某种顺序逐一拆除所有承重墙。拆除第 $i$ 面墙时，产生的拆除能量等于：
\[
\texttt{walls}[\text{左邻}] \times \texttt{walls}[i] \times \texttt{walls}[\text{右邻}]
\]
其中"左邻"和"右邻"是指当前仍存在的、与第 $i$ 面墙相邻的墙。若某一侧已无墙壁，则视该侧结构系数为 1。

拆除一面墙后，该墙消失，其原来的左右邻墙变为直接相邻。

请确定一种拆除顺序，使总拆除能量最大化，返回最大总能量值。

\textbf{示例：}

输入：\texttt{walls = [3, 1, 5, 8]}

输出：\texttt{167}

解释：一种最优顺序为 walls[1] $\to$ walls[2] $\to$ walls[0] $\to$ walls[3]。
拆 walls[1]: $3 \times 1 \times 5 = 15$，剩余 [3,5,8]；
拆 walls[2]: $3 \times 5 \times 8 = 120$，剩余 [3,8]；
拆 walls[0]: $1 \times 3 \times 8 = 24$，剩余 [8]；
拆 walls[3]: $1 \times 8 \times 1 = 8$，剩余 []。
总能量 $= 15+120+24+8 = 167$。

\textbf{提示：}
\begin{itemize}[leftmargin=2em]
  \item $1 \leq n \leq 300$
  \item $0 \leq \texttt{walls}[i] \leq 100$
\end{itemize}



% ╔════════════════════════════════════════════════════════════╗
\category{三、递归 / 记忆化搜索}
% ╚════════════════════════════════════════════════════════════╝

\section{量子密文分裂重组验证}

\noindent\textbf{题目：}

在量子加密通信中，一段明文经过「量子分裂重组」操作后生成密文。操作过程递归定义如下：

\begin{enumerate}[leftmargin=2em]
  \item 若字符串长度为 1，则原样保留。
  \item 若字符串长度 $> 1$，执行以下步骤：
  \begin{enumerate}[label=(\roman*)]
    \item 选择一个分裂位点 $i$（$1 \leq i < $ 字符串长度），将字符串切割为两段：前 $i$ 个字符和后面的字符。
    \item 可选地交换两段的位置（交换或不交换均可）。
    \item 对两段分别递归执行量子分裂重组操作。
    \item 将处理后的两段按序拼接，得到结果。
  \end{enumerate}
\end{enumerate}

给定一段明文 \texttt{plaintext} 和一段密文 \texttt{ciphertext}，判断 \texttt{ciphertext} 能否通过对 \texttt{plaintext} 的某种量子分裂重组操作序列得到。

\textbf{示例 1：}

输入：\texttt{plaintext = "great", ciphertext = "rgeat"}

输出：\texttt{True}

\textbf{示例 2：}

输入：\texttt{plaintext = "abcde", ciphertext = "caebd"}

输出：\texttt{False}

\textbf{提示：}
\begin{itemize}[leftmargin=2em]
  \item $|\texttt{plaintext}| = |\texttt{ciphertext}|$
  \item $1 \leq |\texttt{plaintext}| \leq 30$
  \item 两者均只含小写字母
\end{itemize}



\section{地热勘测最长升温路径}

\noindent\textbf{题目：}

地质勘测队在一片 $m \times n$ 的区域中布设了网格化温度传感器，每个采样点记录了地下恒温层的温度值，形成整数矩阵 \texttt{heatmap}。

勘测队需要绘制一条最长升温路径：从网格中任意一个采样点出发，每一步只能移动到上、下、左、右四个方向的相邻采样点，且移动后的温度必须严格高于当前位置的温度。

请计算这条最长升温路径所经过的采样点数量。

\textbf{示例 1：}

输入：\texttt{heatmap = [[9,9,4],[6,6,8],[2,1,1]]}

输出：\texttt{4}

解释：最长升温路径为 $1 \to 2 \to 6 \to 9$，共 4 个采样点。

\textbf{示例 2：}

输入：\texttt{heatmap = [[3,4,5],[3,2,6],[2,2,1]]}

输出：\texttt{4}

解释：最长升温路径为 $3 \to 4 \to 5 \to 6$，共 4 个采样点。

\textbf{提示：}
\begin{itemize}[leftmargin=2em]
  \item $1 \leq m, n \leq 200$
  \item $0 \leq \texttt{heatmap}[i][j] \leq 2^{31}-1$
\end{itemize}



% ╔════════════════════════════════════════════════════════════╗
\category{四、贪心 / 排序}
% ╚════════════════════════════════════════════════════════════╝

\section{基因序列排序的最小解锁操作}

\noindent\textbf{题目：}

\textbf{题目背景}

在生物研究中，某些基因在不同物种中表现出不同的特征。我们有一个基因序列表示不同物种的基因类型，每个基因有三个可能的类型（1、2、3）。这些基因的排列关系可以影响物种的演化进程。

在这个问题中，基因序列 \texttt{nums} 中的每个元素代表一个基因类型，\texttt{locked} 数组表示基因是否激活。如果 \texttt{nums[i] - nums[i + 1] == 1} 且 \texttt{locked[i] == 0}，则可以交换基因 $i$ 和 $i + 1$（表示基因可以互换位置以促进物种演化）。

我们可以通过交换相邻的基因来尝试使基因序列变得升序排列。每次操作可以解锁一个基因，将 \texttt{locked[i]} 设为 0，从而允许该基因交换位置。

\textbf{要求}

返回将基因序列变成升序所需的最小解锁操作次数。如果不能将基因序列排序，则返回 $-1$。

\textbf{示例 1:}

输入：\texttt{nums = [1, 2, 1, 2, 3, 2], locked = [1, 0, 1, 1, 0, 1]}

输出：\texttt{0}

解释：我们可以按如下方式交换基因序列。交换下标 1 和 2；交换下标 4 和 5。因此，不需要解锁任何基因。

\textbf{示例 2:}

输入：\texttt{nums = [1, 2, 1, 1, 3, 2, 2], locked = [1, 0, 1, 1, 0, 1, 0]}

输出：\texttt{2}

解释：通过解锁下标 2 和 5，我们可以进行如下交换：交换下标 1 和 2；交换下标 2 和 3；交换下标 4 和 5；交换下标 5 和 6。

\textbf{示例 3:}

输入：\texttt{nums = [1, 2, 1, 2, 3, 2, 1], locked = [0, 0, 0, 0, 0, 0, 0]}

输出：\texttt{-1}

解释：尽管所有基因都可以解锁，但无法使基因序列升序排列。

\textbf{提示：}
\begin{itemize}[leftmargin=2em]
  \item 每个基因表示一个不同的特征（1、2、3）。
  \item 通过解锁基因并交换相邻基因位置，模拟基因在物种进化中的变异过程。
\end{itemize}



% ╔════════════════════════════════════════════════════════════╗
\category{五、构造}
% ╚════════════════════════════════════════════════════════════╝

\section{算力网络激活方案（沙地治理）}

\noindent\textbf{题目：}

\textbf{题目描述}

某云计算中心部署了一组新型的三角形算力阵列。该阵列由 \texttt{size} 层计算节点组成，整体呈正三角形分布。为了便于管理，我们将阵列划分为若干个单位计算节点，其中第 $i$ 层（$1 \leq i \leq \texttt{size}$）包含 $2i - 1$ 个计算节点。

初始状态下，所有节点均处于休眠状态。为了节省启动能耗，系统采用了一种"邻域冗余唤醒协议"。你需要手动启动一部分节点作为种子节点，剩余节点将根据以下规则自动唤醒：

\textbf{邻接关系}：若两个计算节点在几何结构上共用一条物理边，则视为相邻。例如，坐标 $(1,1)$ 与 $(2,2)$ 相邻，但 $(2,2)$ 与 $(3,3)$ 不相邻（因没有公共边）。

\textbf{唤醒规则}：若一个休眠节点的相邻节点中，至少有 2 个节点处于工作状态（无论是被手动启动还是被自动唤醒的），该休眠节点将因接收到足够的握手信号而自动切换为工作状态。

\textbf{连锁反应}：被自动唤醒的节点可以继续作为信号源，去唤醒其他相邻的休眠节点。

现需要将一个层数为 \texttt{size} 的算力阵列全部激活（即所有节点转化为工作状态）。请计算并返回最少需要手动启动的节点数量及坐标列表。你可以返回任意一种满足条件的方案。

注意：坐标 $(r, c)$ 表示第 $r$ 层的第 $c$ 个节点，下标从 1 开始。

\textbf{示例 1}

输入：\texttt{size = 3}

输出：\texttt{[[1,1], [2,1], [2,3], [3,1], [3,5]]}

解释：一种方案为手动启动上述 5 个计算节点：
\begin{enumerate}[leftmargin=2em]
  \item 节点 $(2, 2)$ 的相邻节点中，$(1,1)$ 和 $(2,1)$ 已启动，满足 2 个相邻工作节点的条件，因此 $(2,2)$ 被唤醒。
  \item 随后，节点 $(3,2)$ 检测到 $(2,1)$ 和 $(2,2)$ 处于工作状态，也随之唤醒。
  \item 连锁反应持续，最终整个阵列被点亮。
\end{enumerate}

\textbf{示例 2}

输入：\texttt{size = 2}

输出：\texttt{[[1,1], [2,1], [2,3]]}

解释：手动启动 3 个角落的节点。中间的节点 $(2,2)$ 与这三个节点均相邻，满足唤醒条件，从而使整个阵列激活。

\textbf{提示：}
\begin{itemize}[leftmargin=2em]
  \item $1 \leq \texttt{size} \leq 1000$
  \item 返回的坐标必须在有效的阵列范围内，即 $1 \leq r \leq \texttt{size}$ 且 $1 \leq c \leq 2r - 1$。
\end{itemize}



% ╔════════════════════════════════════════════════════════════╗
\category{六、栈 / 单调栈}
% ╚════════════════════════════════════════════════════════════╝

\section{无人机最大矩形着陆区}

\noindent\textbf{题目：}

无人机调度中心正在一个 $m \times n$ 的城市网格中规划应急着陆区。网格中每个格子的状态为：
\begin{itemize}[leftmargin=2em]
  \item \texttt{'1'} —— 空旷区：地面平整，可用于着陆。
  \item \texttt{'0'} —— 障碍区：有建筑物或其他障碍，不可着陆。
\end{itemize}

着陆区必须是一个由连续空旷区格子组成的矩形（各边平行于网格边界）。请计算该城市网格中最大矩形着陆区的面积。

\textbf{示例：}

输入：
\texttt{cityGrid = [["1","0","1","0","0"], ["1","0","1","1","1"], ["1","1","1","1","1"], ["1","0","0","1","0"]]}

输出：\texttt{6}

\textbf{提示：}
\begin{itemize}[leftmargin=2em]
  \item $1 \leq m, n \leq 200$
  \item \texttt{cityGrid[i][j]} 为 \texttt{'0'} 或 \texttt{'1'}
\end{itemize}



% ╔════════════════════════════════════════════════════════════╗
\category{七、单调队列}
% ╚════════════════════════════════════════════════════════════╝

\section{地震监测滑动窗口峰值}

\noindent\textbf{题目：}

地震监测站每秒采集一次地表振动加速度值（单位：gal），形成长度为 $n$ 的时间序列 \texttt{seismicData}。

预警系统需要维护一个大小固定为 $k$ 秒的滑动检测窗口，窗口从时间序列的起始位置开始，每次向后滑动 1 秒。对每个窗口位置，系统需要即时报告窗口内的最大振动加速度（峰值），以判断是否触发预警。

请返回一个数组，包含每个窗口位置的峰值。

\textbf{示例：}

输入：\texttt{seismicData = [1,3,-1,-3,5,3,6,7], k = 3}

输出：\texttt{[3,3,5,5,6,7]}

解释：窗口 $[1,3,-1]$ 峰值=3；窗口 $[3,-1,-3]$ 峰值=3；窗口 $[-1,-3,5]$ 峰值=5；窗口 $[-3,5,3]$ 峰值=5；窗口 $[5,3,6]$ 峰值=6；窗口 $[3,6,7]$ 峰值=7。

\textbf{提示：}
\begin{itemize}[leftmargin=2em]
  \item $1 \leq n \leq 10^5$，$1 \leq k \leq n$
  \item $-10^4 \leq \texttt{seismicData}[i] \leq 10^4$
\end{itemize}



% ╔════════════════════════════════════════════════════════════╗
\category{八、堆 / 数据结构设计}
% ╚════════════════════════════════════════════════════════════╝

\section{实时竞拍中位价引擎}

\noindent\textbf{题目：}

某在线艺术品拍卖平台需要一个高效的价格统计引擎。随着拍卖进行，竞价者不断提交出价。平台需要在任意时刻快速计算当前所有已提交出价的中位数，以便实时调整展示策略和推荐底价。

请设计一个数据结构 \texttt{AuctionMedian}，支持以下操作：
\begin{itemize}[leftmargin=2em]
  \item \texttt{AuctionMedian()} —— 初始化引擎。
  \item \texttt{recordBid(price)} —— 记录一个新的出价 \texttt{price}（整数）。
  \item \texttt{getMedianBid()} —— 返回当前所有已记录出价的中位数（浮点数）。若出价总数为偶数，返回中间两个数的平均值。
\end{itemize}

\textbf{示例：}

\texttt{engine = AuctionMedian()}

\texttt{engine.recordBid(1)}

\texttt{engine.recordBid(2)}

\texttt{engine.getMedianBid()} \quad 返回 1.5

\texttt{engine.recordBid(3)}

\texttt{engine.getMedianBid()} \quad 返回 2.0

\textbf{提示：}
\begin{itemize}[leftmargin=2em]
  \item $-10^5 \leq \texttt{price} \leq 10^5$
  \item 在调用 \texttt{getMedianBid} 前至少已调用过一次 \texttt{recordBid}
  \item 总操作次数不超过 $5 \times 10^4$
\end{itemize}



% ╔════════════════════════════════════════════════════════════╗
\category{九、回溯}
% ╚════════════════════════════════════════════════════════════╝

\section{N座激光基站无干扰部署}

\noindent\textbf{题目：}

国防部需要在一个 $n \times n$ 的战术网格中部署 $n$ 座激光通信基站。每座基站一旦启动，将同时向以下四个方向发射激光束：
\begin{itemize}[leftmargin=2em]
  \item 所在行的所有格子
  \item 所在列的所有格子
  \item 所在主对角线方向（左上 $\leftrightarrow$ 右下）的所有格子
  \item 所在副对角线方向（右上 $\leftrightarrow$ 左下）的所有格子
\end{itemize}

为了避免激光干扰，任意两座基站不能处于同一行、同一列或同一条对角线上。

请返回所有合法的部署方案。每个方案用一个字符串数组表示，其中 \texttt{'Q'} 代表基站位置，\texttt{'.'} 代表空位。

\textbf{示例：}

输入：\texttt{n = 4}

输出：\texttt{[[".Q..","...Q","Q...","..Q."], ["..Q.","Q...","...Q",".Q.."]]}

\textbf{提示：}
\begin{itemize}[leftmargin=2em]
  \item $1 \leq n \leq 9$
\end{itemize}



% ╔════════════════════════════════════════════════════════════╗
\category{十、二分查找}
% ╚════════════════════════════════════════════════════════════╝

\section{深空信号分段传输}

\noindent\textbf{题目：}

深空探测器在一次远征任务中采集了 $n$ 个非负整数组成的观测数据序列 \texttt{payload}。由于星际通信信道的带宽限制，这些数据必须被分成恰好 $k$ 个非空的\textbf{连续}片段，分批传回地球。

每个片段的传输代价等于该片段内所有数据之和。由于通信协议的峰值带宽约束，传输代价最大的那个片段将成为整个传输过程的瓶颈。

请设计一种分段方案，使得所有片段中传输代价的\textbf{最大值}尽可能小。返回该最小化后的最大传输代价。

\textbf{示例 1：}

输入：\texttt{payload = [7,2,5,10,8], k = 2}

输出：\texttt{18}

解释：最优分段为 $[7,2,5]$ 和 $[10,8]$，两段代价分别为 14 和 18，瓶颈代价为 18。

\textbf{示例 2：}

输入：\texttt{payload = [1,2,3,4,5], k = 2}

输出：\texttt{9}

解释：最优分段为 $[1,2,3]$ 和 $[4,5]$，两段代价分别为 6 和 9。

\textbf{提示：}
\begin{itemize}[leftmargin=2em]
  \item $1 \leq n \leq 1000$
  \item $0 \leq \texttt{payload}[i] \leq 10^6$
  \item $1 \leq k \leq \min(50, n)$
\end{itemize}



% ╔════════════════════════════════════════════════════════════╗
\category{十一、图论}
% ╚════════════════════════════════════════════════════════════╝

\section{骨干光缆抗毁性分析}

\noindent\textbf{题目：}

某电信运营商的骨干通信网络由 $n$ 个数据中心（编号 $0$ 到 $n-1$）和若干条双向光缆组成。网络中任意两个数据中心之间当前都可以互相通信（直接或间接）。

运营商的安全团队需要进行\textbf{抗毁性分析}：找出所有"关键光缆"——即一旦该光缆被切断（例如被自然灾害破坏），将导致至少两个数据中心之间\textbf{无法}互相通信。

给定数据中心数量 $n$ 和光缆连接列表 \texttt{cables}（每条光缆表示为 $[a, b]$，表示数据中心 $a$ 和 $b$ 之间有一条双向光缆），返回所有关键光缆的列表。

\textbf{示例：}

输入：\texttt{n = 4, cables = [[0,1],[1,2],[2,0],[1,3]]}

输出：\texttt{[[1,3]]}

解释：光缆 $[1,3]$ 是唯一的关键光缆。切断它后，数据中心 3 将与其余中心断联。

\textbf{提示：}
\begin{itemize}[leftmargin=2em]
  \item $2 \leq n \leq 10^5$
  \item $n - 1 \leq |\texttt{cables}| \leq 10^5$
  \item $0 \leq a_i, b_i \leq n - 1$，$a_i \neq b_i$
  \item 不存在重复光缆，初始时所有数据中心互相可达
\end{itemize}



% ╔════════════════════════════════════════════════════════════╗
\category{十二、链表}
% ╚════════════════════════════════════════════════════════════╝

\section{数据包分段逆序传输}

\noindent\textbf{题目：}

在一个受限的卫星通信协议中，数据以单链表形式排列。为了增强传输过程中的纠错能力，发送端需要对数据链进行\textbf{分段逆序编码}：将链中的节点每 $k$ 个为一组，对每组内部进行逆序排列，但组与组之间的相对顺序不变。

如果链末尾剩余节点不足 $k$ 个，则这些节点\textbf{保持原始顺序}不做处理。

编码过程只能改变节点之间的链接关系，\textbf{不能修改}节点内部存储的数值。

\textbf{示例 1：}

输入：\texttt{dataFrame = [1,2,3,4,5], k = 2}

输出：\texttt{[2,1,4,3,5]}

\textbf{示例 2：}

输入：\texttt{dataFrame = [1,2,3,4,5], k = 3}

输出：\texttt{[3,2,1,4,5]}

\textbf{提示：}
\begin{itemize}[leftmargin=2em]
  \item 链中节点数 $n$ 满足 $1 \leq n \leq 5000$
  \item $0 \leq$ 节点值 $\leq 1000$
  \item $1 \leq k \leq n$
\end{itemize}


\end{document}
